\chapter{Problem Formulation}
\label{chap:foundations}
This chapter introduces the shared formal setup used throughout the
thesis. We first formulate the goal-conditioned RL problem
(\S\ref{sec:problem}) and then describe the latent action-conditioned
world model assumed by the later methods (\S\ref{sec:lewm-bg}). The
description is kept independent of a particular implementation: the
LeWM instantiation used in the experiments is described in
\S\ref{sec:exp-wm-instantiation}, and its predictive accuracy, latent
geometry, and cross-environment variation are evaluated in
\S\S\ref{sec:wm-quality}--\ref{sec:wm-quality-envs}.

% ---------------------------------------------------------------------
\section{Goal-Conditioned Control in Latent Space}
\label{sec:problem}

We formulate robotic goal-reaching as a goal-conditioned Markov
decision process (GCMDP)

\begin{equation}
\mathcal{M}_g =
(\mathcal{S}, \mathcal{A}, \mathcal{G}, p, r, \gamma),
\end{equation}
where $\mathcal{S}$ is the observation space,
$\mathcal{A}=\mathbb{R}^{d_a}$ is a continuous action space,
$\mathcal{G}\subseteq\mathcal{S}$ denotes the set of valid goals,
$p(s'|s,a)$ describes the environment dynamics,
$r(s,a,g)$ is a goal-conditioned reward function, and
$\gamma\in[0,1)$ is the discount factor. The objective is to learn a
goal-conditioned policy $\pi(a|s,g)$ that reaches arbitrary target
states $g$.

In the environments considered throughout this thesis, observations are
high-dimensional RGB images. Learning and planning directly from pixels
is computationally expensive and sample-inefficient, particularly over
long horizons.
Instead, all methods operate in a learned latent space

\begin{equation}
\mathcal{Z}=\mathbb{R}^{d_z},
\end{equation}
obtained through an encoder
$E_\xi:\mathcal{S}\rightarrow\mathcal{Z}$.
The resulting latent-space GCMDP is

\begin{equation}
\widehat{\mathcal{M}}_g
=
(\mathcal{Z}, \mathcal{A}, \mathcal{Z}, \hat p, \hat r, \gamma),
\end{equation}
where both the current state and the goal are represented as latent
vectors.

State transitions in this latent space are modelled by an
action-conditioned predictor

\begin{equation}
\hat p(z'|z,a) \approx f_\psi(z,a),
\end{equation}
which predicts the next latent state given the current latent and
action. Together, the encoder and predictor,

\begin{equation}
(E_\xi,f_\psi),
\end{equation}
form the learned world model described in
\S\ref{sec:lewm-bg}. Once trained, the world model is frozen and reused
by all downstream algorithms, providing both a compact state
representation and a differentiable simulator for imagined rollouts.

Operating in latent space avoids modelling pixels directly and gives
the RL algorithms compact task-relevant features. Because the space is
regularised to be approximately isotropic (\S\ref{sec:lewm-bg}),
Euclidean distance can also be used for simple rewards and goal metrics.

% ---------------------------------------------------------------------
\section{Latent Action-Conditioned World Model}
\label{sec:lewm-bg}

The methods in this thesis use a latent action-conditioned world model
with two components: an encoder $E_\xi$ that maps observations to latent
states, and a latent dynamics predictor $f_\psi$ that advances those
states under actions. This section describes the assumptions needed by
the downstream algorithms. The concrete implementation used in the
experiments, LeWM \cite{maes2026leworldmodelstableendtoendjointembedding},
is described in \S\ref{sec:exp-wm-instantiation}.

At a high level, the model follows the Joint-Embedding
Predictive Architecture (JEPA) paradigm
(\S\ref{sec:jepa}, Figure~\ref{fig:jepa-bg}). Rather
than reconstructing future images, the model learns a latent
representation in which future states can be predicted directly. This
produces a compact state space that captures the aspects of the
environment most relevant to control while avoiding the complexity of
pixel-level prediction.

\paragraph{Structure.}

The encoder maps each observation $s_t$ to a latent state

\begin{equation}
z_t = E_\xi(s_t) \in \mathbb{R}^{d_z}.
\end{equation}
Actions are embedded into the same latent space and provided as inputs
to the dynamics model. The predictor $f_\psi$ receives a short history
of latent states together with an action chunk and predicts the
corresponding future latent state. Throughout this thesis, a single
world-model transition corresponds to one chunk of $h$ primitive
environment actions, meaning that one prediction step advances the
latent state by $h$ simulator frames.

The resulting dynamics model can therefore be viewed as a learned
transition function operating directly in latent space,

\begin{equation}
z_{t+1} = f_\psi(z_t,a_t),
\end{equation}
where, for brevity, the action chunk notation is omitted.

\paragraph{Predictive training objective.}

The world model is trained using a self-supervised prediction
objective. Given a sequence of observations and actions, the predictor
is required to match the latent representation of a future state:

\begin{equation}
\mathcal{L}_{\mathrm{pred}}
=
\mathbb{E}_{\tau\sim\mathcal D}
\left[
\big\|
f_\psi(z_{t-c},a_{t-c:t-1})
-
z_t
\big\|_2^2
\right],
\label{eq:pred-loss-clean}
\end{equation}
where $z_t=E_\xi(s_t)$.

Unlike pixel reconstruction objectives, this loss encourages the model
to retain only the information required for predicting future
dynamics. The learned latent space therefore becomes a compact,
control-oriented representation of the environment.

\paragraph{Isotropy regularisation.}

A challenge in JEPA-style training is avoiding representation collapse,
where all observations are mapped to nearly identical latent vectors.
To prevent this, the model includes an isotropy regulariser that keeps
the latent distribution spread out and approximately Gaussian.

The complete training objective is

\begin{equation}
\mathcal{L}_{\mathrm{WM}}
=
\mathcal{L}_{\mathrm{pred}}
+
\lambda_{\mathrm{iso}}
\mathcal{L}_{\mathrm{iso}}.
\label{eq:lewm-total-clean}
\end{equation}
Here $\mathcal{L}_{\mathrm{iso}}$ denotes the isotropy penalty used by
the chosen implementation; in the experiments this is SIGReg.

\paragraph{Usage in this work.}
The encoder provides the latent state representation used throughout
the thesis, while the predictor acts as a differentiable simulator for
imagined rollouts. All policies, value functions, rewards, and planners
operate exclusively on latent states
$z\in\mathcal Z$ and predicted transitions generated by
$f_\psi$ unless otherwise stated.

\paragraph{Implications for reinforcement learning and planning.}

Two properties matter for the methods developed in subsequent chapters.
First, the dynamics predictor enables fast latent-space simulation,
allowing candidate actions and subgoals to be evaluated without
expensive physics rollouts.
Second, the isotropy regularisation induces a structured latent
geometry in which Euclidean distance is a meaningful measure of state
similarity. Empirically,
states corresponding to similar configurations of the environment are
embedded close together, while dissimilar states are embedded further
apart. As a result, the distance

\begin{equation}
\|z-g\|_2
\end{equation}
provides a useful proxy for progress towards a goal state $g$. This
proxy is useful but not, on its own, a sufficient control objective: as
\S\ref{sec:exp-cem-baseline} later shows, planning to minimise latent
distance alone reaches the goal \emph{embedding} while failing the actual
task.
This geometry is used throughout the thesis: it underpins the
latent-space rewards in Chapter~\ref{chap:wm-env} and forms part of the
value-based scoring functions used by the planning methods in
Chapter~\ref{chap:wm-planner}.

% ---------------------------------------------------------------------
\section{Reward Functions in Latent Space}
\label{sec:wmep-setup}

When a method treats the world model as a simulated environment, rewards
are computed directly in latent space using one of two reward functions.
For sparse-reward experiments, success is defined by reaching a latent
state within a threshold distance $\delta$ of the goal,

\begin{equation}
\tilde r_t^{\mathrm{sparse}}
=
\mathbb{1}\!\left[\|z_t-z_g\|_2 < \delta\right].
\end{equation}
For dense-reward experiments, the reward is the negative distance to
the goal,

\begin{equation}
\tilde r_t^{\mathrm{dense}}
=
-\|z_t-z_g\|_2.
\end{equation}
The dense reward provides a shaped learning signal at every timestep,
whereas the sparse reward only indicates whether the goal has been
reached. To keep offline and online data consistent, trajectories in the
offline replay buffer are relabelled using the same reward function
whenever dense rewards are employed. This avoids training the critic on
a mixture of reward conventions.
